\documentclass[11pt,a4paper]{article}
\usepackage[T1]{fontenc}
\usepackage[utf8]{inputenc}
\usepackage{lmodern}
\usepackage[a4paper,margin=27mm,headheight=14pt]{geometry}
\usepackage{amsmath,amssymb,amsthm,mathtools}
\usepackage{booktabs,tabularx,array}
\usepackage{microtype}
\usepackage{xcolor}
\usepackage{enumitem}
\usepackage{fancyhdr}
\usepackage[hidelinks]{hyperref}
\usepackage{bookmark}
\usepackage{xurl}
\definecolor{ink}{HTML}{17384A}
\definecolor{accent}{HTML}{256477}
\hypersetup{pdftitle={Constructive Diophantine Theorems in iPRA},
  pdfauthor={Research report prepared for Vladimir Reshetnikov},
  pdfsubject={Størmer, quantitative Davenport--Schmidt, and the effective boundary for Roth}}
\newtheorem{theorem}{Theorem}[section]
\newtheorem{lemma}[theorem]{Lemma}
\newtheorem{proposition}[theorem]{Proposition}
\newtheorem{corollary}[theorem]{Corollary}
\theoremstyle{definition}
\newtheorem{definition}[theorem]{Definition}
\theoremstyle{remark}
\newtheorem{remark}[theorem]{Remark}
\newcommand{\ipra}{\ensuremath{\mathrm{iPRA}}}
\newcommand{\Q}{\mathbb Q}
\newcommand{\Z}{\mathbb Z}
\newcommand{\N}{\mathbb N}
\setlist{itemsep=4pt,topsep=6pt}
\setlength{\parskip}{4pt plus 1pt}
\setlength{\parindent}{0pt}
\emergencystretch=2em
\pagestyle{fancy}
\fancyhf{}
\fancyhead[L]{\small\color{ink}Constructive Diophantine theorems in iPRA}
\fancyhead[R]{\small 15 November 2026}
\fancyfoot[C]{\thepage}
\renewcommand{\headrulewidth}{0.3pt}

\begin{document}
\begin{titlepage}
\vspace*{16mm}
{\large\color{accent}\sffamily RESEARCH REPORT}\par
\vspace{8mm}
{\Huge\color{ink}\bfseries Constructive Diophantine\\[3pt]
Theorems in iPRA}\par
\vspace{9mm}
{\Large St\o rmer, quantitative Davenport--Schmidt,\\[4pt]
and the effective boundary for Roth}\par
\vspace{12mm}
{\large Prepared for Vladimir Reshetnikov}\par
{\large 15 November 2026}\par
\vspace{15mm}
\hrule
\vspace{5mm}
\textbf{What is proved.}
A complete finite enumeration theorem for consecutive smooth integers;
a complete exponent-three approximation theorem with explicit primitive
recursive bounds and positive separation data; and a complete effective
Roth theorem for quadratic algebraic inputs, with a general Liouville bound.
\vspace{5mm}

\textbf{What remains unresolved.}
The general effective Roth cutoff is not proved.  It would resolve a major
open effectivity problem.  The new Lean files are specifications and a
formalization route, not completed machine-checked iPRA derivations.
\vspace{5mm}
\hrule
\vfill
{\small Based on ProveIt commit
\texttt{b33ee1acb02b7619dfef419e61c3e343735a7bc4}.\\
The archive contains this article, its LaTeX source, exact reference
computations, verification evidence, and Lean specifications.}
\end{titlepage}

\begin{abstract}
We develop finite constructive versions of three Diophantine statements
in the setting of first-order intuitionistic primitive recursive arithmetic.
For St\o rmer's theorem, a cubic-free coefficient reduction and an elementary
Pell argument give a complete primitive recursive enumeration of all
consecutive smooth pairs, together with explicit cardinality and size bounds.
For Davenport--Schmidt approximation, we reconstruct the minimal-point proof
with explicit finite search bounds and a single rational sample.  Given a
primitive recursive real name and a primitive recursive quadratic separation
modulus, the construction produces arbitrarily high rational or quadratic
approximants with error less than
$32768\max(1,\xi^2)H^{-3}$.  Resultants provide the required modulus for
algebraic inputs of degree at least three.  For Roth, we explain why a general
witness-producing cutoff would solve the unresolved effective Roth problem,
and prove instead the full quadratic case and the all-degree Liouville range.
Every positive result is accompanied by a finite arithmetic proof and a
specific path to object-level Lean derivations in ConciseIPRA.  Those new
derivations are not claimed to be completed.
\end{abstract}
\tableofcontents
\clearpage
\section{Scope, results, and the arithmetic framework}
\label{sec:framework}

The three requested theorems do not have the same constructive status.  A
finite list of smooth neighbors can be computed; the existence of good
quadratic approximants can be made quantitative when nonquadraticity is
supplied as positive separation data; and a general denominator cutoff in
Roth's theorem remains an outstanding effectivity problem.  These distinctions
are mathematical requirements, rather than limitations of a particular proof
assistant.  We therefore state precisely what has been established.

\begin{center}
\begin{tabularx}{\textwidth}{@{}p{0.19\textwidth}X@{}}
\toprule
Requested result & Result established in this report\\
\midrule
St\o rmer & A complete primitive recursive enumeration, with at most
$3^k-2^k$ neighbors and an explicit bound on every entry.\\[3pt]
Davenport--Schmidt & Exponent $3$, a larger absolute constant, and an
explicit height search bound, for a fixed primitive recursive real name
with a fixed primitive recursive quadratic separation modulus.
Algebraic inputs of degree at least three have an explicit such modulus.\\[3pt]
Roth & A complete cutoff proof for every quadratic algebraic irrational
and every positive rational exponent increment; a complete all-degree
Liouville bound; an exact account of the unresolved general cutoff.\\[3pt]
Lean & Concrete arithmetic specifications, an actual repository transport
interface, and a dependency-by-dependency plan.  The new Diophantine
derivations have not been completed or kernel checked.\\
\bottomrule
\end{tabularx}
\end{center}

In particular, this is \emph{not a proof of all three requested full theorems}.
Calling a number-of-solutions estimate an enumeration algorithm would conceal
the central difficulty in Roth's theorem.  Section~\ref{sec:roth} explains the
distinction and proves the fragment that is actually claimed.

\subsection{The inspected object theory}

We use the first-order intuitionistic primitive recursive arithmetic defined
in \path{ConciseIPRA/PR.lean}, at ProveIt commit
\begin{center}
\texttt{b33ee1acb02b7619dfef419e61c3e343735a7bc4}.
\end{center}
The inspected source defines the primitive recursive function symbols
\texttt{PR} by zero, successor, projections, composition, and primitive
recursion.  Its terms are \texttt{PTm}; its axioms \texttt{PRAAx} are their
defining equations, successor nonzero and injectivity, and induction for
quantifier-free formulas.  The underlying derivations are the intuitionistic
natural-deduction relation \texttt{Prf} in \texttt{Deduction.lean}, with
equality and all first-order connectives and quantifiers
\cite{ProveItPR,ProveItDeduction}.

Thus ``in iPRA'' has a specific meaning.  A final formalization must construct
an inhabitant of
\[
 \texttt{Prf PRAAx [] statement},
\]
where \texttt{statement} is an encoded first-order formula.  A Lean theorem
about the standard natural numbers or real numbers is a useful intermediate
result, but is not that object-theory derivation.  The inspected
\texttt{BiInterpretation.lean} proves \texttt{pra\_theorem\_iff}, which transfers
such a derivation to the concise scalar theory by its existing translation
\cite{ProveItBi}.  This transfer does not create a missing arithmetic proof.

We also inspected the Dirichlet--Kronecker and Minkowski archives in
\texttt{docs/drafts} at the same commit.  Their finite arithmetic perspective
is relevant; their own status files identify their approximation artifacts as
specifications or unfinished formalization work.  No theorem from those
archives is assumed here.  In particular, the two pigeonhole constructions
and every Diophantine estimate needed below are proved in this report.

\subsection{Finite codes and quantifier-free induction}

Signed integers are pairs of natural numbers, or a sign bit and a magnitude;
rationals are signed numerators with positive denominators.  Finite vectors,
polynomials, and lists are encoded by elementary pairing and list operations.
The choice of codec has no bearing on the estimates.  Malformed codes receive
fixed default outputs; correctness statements include a decidable input
validity hypothesis.

Addition, multiplication, natural powers, quotient and remainder, absolute
value, gcd, integer square root, factorials, bounded sums and products,
and bounded minimization have primitive recursive definitions.  A bounded
universal condition is a finite Boolean conjunction and a bounded existential
condition a finite Boolean disjunction.  In this expanded language, each is
represented by an equation between primitive recursive terms.  We use
quantifier-free induction on these characteristic equations, with numerical
parameters allowed.

For example, finite pigeonhole is not an infinite combinatorial principle.
Given $N+1$ colors in $\{0,\ldots,N-1\}$, scan the list and maintain a finite
occupancy table.  Either a repeated color is returned with its two indices,
or each new item increases the number of occupied entries by one.  That number
is at most $N$, so the latter branch is impossible after $N+1$ steps.  The
table invariant and count identity are bounded Boolean assertions.  The
algorithm's final success flag is decidable; eliminating its false branch is
intuitionistically legitimate.

The same discipline applies to a ``least'' object.  We only select least
elements of explicitly bounded finite collections.  No unrestricted
least-number principle, unbounded minimization, compactness theorem, or
choice of an infinite subsequence is used in the positive results.

\subsection{Real notation and positive input information}

A fast rational name is a fixed primitive recursive definition $r(k)$,
with arithmetic correctness conditions expressing an error at most $2^{-k}$.
Notation such as $|\xi-r(k)|\leq2^{-k}$ abbreviates those conditions; it does
not add a completed real-number sort to first-order iPRA.  Strict inequalities
must have positive rational error allowances.  Closed inequalities are read
through their rational error estimates and used at explicitly chosen finite
precisions.

There are two equivalent ways to organize the resulting mathematics.  One
can state a theorem schema for each fixed choice of primitive recursive name
and modulus.  Alternatively, one can state a single finite theorem about
rational approximations, finite tables, and rational certificates, then apply
it to the fixed name.  The Davenport--Schmidt proof makes this second route
explicit by reducing its argument to a single rational approximation at a
specified precision.  This matters because induction on arbitrary assertions
about real limits would exceed the stated induction scheme.

There is no total primitive recursive evaluator for every primitive recursive
program supplied as unrestricted input.  The schemas used here neither
postulate nor need such an evaluator.  Polynomial evaluation and finite-table
checking, by contrast, are particular fixed primitive recursive algorithms.

For St\o rmer's theorem even a real-name schema is unnecessary: all radicals
are elements of an explicitly coded ordered quadratic field, with sign
decided by integer squaring.  For Roth's proved fragment the algebraic name is
produced by a fixed rational bisection algorithm.  For Davenport--Schmidt the
separation modulus is additional positive input, and its necessity is not
hidden under the phrase ``not rational or quadratic.''

\subsection{Status of the proofs and of their implementation}

The positive results below have complete human-readable finite mathematical
proofs, including explicit bounds.  Their proposed iPRA formalizations use
only the described arithmetic definitions and finite induction invariants.
The remaining Lean work is substantial: it must encode the arithmetic,
establish each invariant in \texttt{Prf PRAAx}, and assemble the derivations.
The included code tests computations and certificates, not the existence of
those derivations.  Neither successful Python examples nor an absence of
\texttt{sorry} in a specification file is presented as a completed
formalization.

\section{Roth: the effective boundary and a complete quadratic theorem}
\label{sec:roth}

\subsection{Three different meanings of finiteness}

For a real algebraic irrational $\alpha$ and rational $\delta=a/b>0$, where
$a,b\geq1$ are integers, Roth's theorem says that the reduced fractions
$p/q$ satisfying
\begin{equation}\label{eq:roth-inequality}
  \left|\alpha-\frac pq\right|<q^{-2-a/b},\qquad q\geq1,
\end{equation}
form a finite set.  A constructive formulation must specify the data meant by
``finite.''  The following are different assertions.
\begin{enumerate}
\item A \emph{cutoff}: produce an integer $Q$ such that every solution has
$q<Q$.  This permits a finite exhaustive search.
\item A \emph{cardinality bound}: produce $N$ and prove that no list of $N+1$
distinct reduced fractions consists entirely of solutions.
\item A \emph{finite obstruction}: for every supplied finite list satisfying
specified growth conditions, locate a member that is not a solution.
\end{enumerate}
A cutoff is not supplied by an upper bound on cardinality.  Indeed, let
$P_e(n)$ say that program $e$ halts for the first time at step $n$.  This is a
primitive recursive predicate, and for every $e$ there is at most one $n$
with $P_e(n)$.  A computable function $B(e)$ bounding all such $n$ would
solve the halting problem by simulating $e$ for $B(e)$ steps.  Thus no general
algorithm converts cardinality bounds into cutoffs.

Roth's theorem is known to have effective cardinality estimates; Luckhardt's
analysis of the Roth and Esnault--Viehweg proofs is a particularly relevant
example \cite{Luckhardt1989,KohlenbachHerbrand}.  The general cutoff problem
is a different, unresolved effectivity problem.  The introduction of
Adiceam and Shirandami's 2026 paper explicitly describes the corresponding
height problem for the Subspace Theorem, including Roth's case, as open
\cite{AdiceamShirandami2026}.  Their probabilistic results do not provide a
cutoff for each specified algebraic number.

Consequently, a uniform witness-producing proof of
\begin{equation}\label{eq:effective-roth-target}
 \forall A,a,b\ \exists Q\ \forall p,q\quad
 \bigl(\operatorname{Valid}(A)\wedge a,b,q\geq1
       \wedge\operatorname{Good}_{A,a,b}(p,q)\bigr)\Longrightarrow q<Q
\end{equation}
would settle the effective Roth problem, if $A$ ranges over finite algebraic
number codes and the premise is arranged in the usual decidable-input form.
Here and below one can equivalently move the input validity premises outside
the existential quantifier.  A constructive proof in the intended arithmetic
would supply a computable bound; under the explicit primitive recursive
interpretation required in this report it must supply a primitive recursive
one.  This observation is a barrier to claiming the requested full theorem,
\emph{not} a proof that no such proof exists.  Nor does it rule out an iPRA
proof of a suitable quantitative finite-obstruction version.  Such a proof
would require a complete arithmetic treatment of Roth's auxiliary-polynomial
and index estimates.  Those estimates are not proved in this report and are
not smuggled into an axiom named ``effective Roth.''

We next prove an unconditional theorem that does have the required explicit
arithmetic content: the full Roth cutoff for quadratic algebraic irrationals,
and the Liouville range for every degree.  We give both the familiar
human-readable argument and its entirely rational, finitely checkable form.

\subsection{Finite input data and the separation constant}

Let
\[
 f(X)=\sum_{i=0}^{d}c_iX^i\in\mathbb Z[X],\qquad d\geq2,\quad c_d\ne0.
\]
Assume $f$ has no rational root.  This assumption admits a finite certificate.
Require $c_0\ne0$ and check
\begin{equation}\label{eq:no-root-certificate}
 \sum_{i=0}^{d}c_i u^i v^{d-i}\ne0
 \quad\text{for every } |u|\leq |c_0|,\quad 1\leq v\leq |c_d|.
\end{equation}
There is no need to restrict the finite check to coprime pairs.  The
rational-root theorem proves its sufficiency: if a reduced $p/q$ is a root,
reduction of $q^df(p/q)=0$ modulo $p$ and modulo $q$, followed by Euclid's
lemma, gives $p\mid c_0$ and $q\mid c_d$.  The case $p=0$ is excluded by
$c_0\ne0$.  Therefore $|p|\leq|c_0|$ and $q\leq|c_d|$, contradicting the
check.  The same conclusion then holds for unreduced fractions by reducing
them.  Euclid's lemma here follows from the integer extended Euclidean
algorithm, or directly from its finite Bezout certificate.

Take rational numbers $u_0<v_0$ with $f(u_0)f(v_0)<0$, and an integer $R\geq1$
with $|u_0|,|v_0|\leq R$.  These are all decidable rational conditions.  Put
\begin{equation}\label{eq:roth-K}
 K=\max\left(1,\sum_{i=1}^{d} i|c_i|(R+1)^{i-1}\right).
\end{equation}
For rational or real $x,y$ with $|x|,|y|\leq R+1$, the identity
\[
 x^i-y^i=(x-y)\sum_{j=0}^{i-1}x^{i-1-j}y^j
\]
gives
\begin{equation}\label{eq:polynomial-lipschitz}
 |f(x)-f(y)|\leq K|x-y|.
\end{equation}
This proof uses only finite sums, the triangle inequality, and multiplication;
no derivative, mean-value theorem, compactness principle, or maximum is
needed.

The input defines a primitive recursive rational name $r_k$.  Starting from
$[u_0,v_0]$, bisect $2R+k$ times, retaining the half whose endpoint values
have opposite signs, and let $r_k$ be the left endpoint.  Every midpoint is
rational and cannot be a root by \eqref{eq:no-root-certificate}, so the
branch is decided by an integer sign computation.  Write $s_k$ for the right
endpoint after the same number of steps.  Since
$v_0-u_0\leq2R\leq2^{2R}$, the construction gives
\begin{equation}\label{eq:root-name-properties}
 \begin{split}
 &|r_k|,|s_k|\leq R,\qquad 0<s_k-r_k\leq2^{-k},\\
 &|r_k-r_l|\leq2^{-k}+2^{-l},\qquad |f(r_k)|\leq K2^{-k}.
 \end{split}
\end{equation}
For the last inequality, opposite signs imply
$|f(r_k)|\leq|f(r_k)-f(s_k)|$, and then
\eqref{eq:polynomial-lipschitz} applies.  For the Cauchy inequality, when
$k\leq l$ the later interval lies inside the earlier one, so the difference
between left endpoints is at most $2^{-k}$; the other order is symmetric.
These are finite induction arguments about a bounded bisection computation.
The symbol $\alpha$ below denotes the real described by this name; the
arithmetic proof itself needs only \eqref{eq:root-name-properties}.

\subsection{Liouville separation without limits}

\begin{theorem}[Explicit algebraic separation]\label{thm:liouville-explicit}
For every integer $p$, positive integer $q$, and $k\geq0$,
\begin{equation}\label{eq:finite-separation}
 \left|\frac pq-r_k\right|+2^{-k}\ \geq\ \frac1{Kq^d}.
\end{equation}
In the conventional real notation this implies
\begin{equation}\label{eq:liouville-real}
 \left|\frac pq-\alpha\right|\geq\frac1{Kq^d}.
\end{equation}
\end{theorem}

\begin{proof}
Let $t=p/q$.  The number
\[
 N_f(p,q)=q^df(p/q)=\sum_{i=0}^{d}c_i p^i q^{d-i}
\]
is a nonzero integer, and therefore
\begin{equation}\label{eq:integer-gap}
 |f(t)|\geq q^{-d}.
\end{equation}
The rational comparison $|p|>(R+1)q$ is decidable.  If it holds, then
$|t-r_k|>1$, while $1/(Kq^d)\leq1$, so the assertion follows.
In the remaining case $|t|\leq R+1$.  For any $l$, equations
\eqref{eq:polynomial-lipschitz}, \eqref{eq:root-name-properties}, and
\eqref{eq:integer-gap} give
\[
 q^{-d}\leq|f(t)|
 \leq |f(t)-f(r_l)|+|f(r_l)|
 \leq K|t-r_l|+K2^{-l}.
\]
Taking $l=k$ and dividing by the positive integer $K$ proves
\eqref{eq:finite-separation} immediately.  Every comparison in this
argument is rational and decidable.
The real formulation is its usual Cauchy-name interpretation.
\end{proof}

For a reader working directly with real algebraic numbers, there is a shorter
presentation of the same argument.  If $|p|\leq(R+1)q$, then $f(\alpha)=0$
and \eqref{eq:polynomial-lipschitz} imply
$q^{-d}\leq|f(p/q)|\leq K|p/q-\alpha|$; the complementary rational branch
again gives an error exceeding $1$.  The expanded proof above explains why
this familiar argument does not conceal an infinitary principle.

\subsection{An explicit Roth cutoff for every quadratic input}

Fractional powers need not appear in the arithmetic language.  Define a
\emph{positive certificate of Roth approximation} to be a triple $(p,q,k)$
with $p\in\mathbb Z$, $q\geq1$, and
\begin{equation}\label{eq:good-certificate}
 q^{2b+a}\left(\left|\frac pq-r_k\right|+2^{-k}\right)^b<1.
\end{equation}
This is a decidable rational inequality, and clearing denominators makes it
an integer inequality.  It asserts a strict upper bound with an explicit
error allowance.  Conversely, a strict real inequality
\eqref{eq:roth-inequality} supplies such a certificate once a positive
rational gap and a sufficient Cauchy precision are given.  No trichotomy
principle for arbitrary real numbers is being assumed.

\begin{theorem}[Primitive recursive cutoff in the Liouville range]
\label{thm:quadratic-roth}
Let $a,b\geq1$ and suppose the integer
\[
 c=a-(d-2)b
\]
is positive.  Every certificate \eqref{eq:good-certificate} satisfies
\begin{equation}\label{eq:cutoff-explicit}
 q^c<K^b,\qquad\text{and hence}\qquad q<K^b.
\end{equation}
In particular, when $d=2$, this proves the full Roth finiteness conclusion
with the explicit cutoff $Q=K^b$ for every positive rational $\delta=a/b$.
\end{theorem}

\begin{proof}
Write $E=|p/q-r_k|+2^{-k}$.  The preceding theorem gives
$E\geq1/(Kq^d)$.  Thus \eqref{eq:good-certificate} implies
\[
 1>q^{2b+a}E^b
   \geq\frac{q^{2b+a}}{K^bq^{db}}
   =\frac{q^c}{K^b}.
\]
As $c\geq1$ and $q\geq1$, finite induction gives $q\leq q^c$.
This proves both inequalities.  For $d=2$, $c=a\geq1$ automatically.
\end{proof}

All approximants are therefore found among the finite pairs
\[
 1\leq q<K^b,\qquad |p|\leq (R+1)q.
\]
The numerator bound follows because a good approximation has error less
than $1$.  Exact membership can be decided from the algebraic input by
sign determination, but even without implementing sign determination the
finite rectangle is already the promised cutoff certificate.  Coprimality
can be checked by the Euclidean algorithm; the proof of the bound did not
require it.

For degree $d\geq3$, the hypothesis is precisely
$\delta>d-2$.  It is essential to identify this as the Liouville range,
rather than to call it general Roth.  The difficult range
$0<\delta\leq d-2$ is not settled by this argument.  For example, the input
$f(X)=X^3-2$ with $0<\delta\leq1$ is outside the theorem's scope.

\paragraph{Concrete example.}
For $f(X)=X^2-2$, take $[u_0,v_0]=[1,2]$ and $R=2$.
The finite rational-root check succeeds, and \eqref{eq:roth-K} gives $K=6$.
Thus every fraction satisfying
$|\sqrt2-p/q|<q^{-2-a/b}$ has $q<6^b$.
For $\delta=1/2$, the explicit cutoff is $q<36$.
These deliberately generous constants make the complete arithmetic proof
transparent; sharper quadratic constants are unnecessary for the theorem.

\subsection{What the elementary gap principle does and does not give}

Suppose two distinct reduced fractions $p/q$ and $r/s$, with $1\leq q\leq s$,
both satisfy \eqref{eq:roth-inequality}.  Since $ps-rq$ is a nonzero integer,
\[
 \frac1{qs}\leq\left|\frac pq-\frac rs\right|
 <q^{-2-\delta}+s^{-2-\delta}\leq2q^{-2-\delta},
\]
and hence
\begin{equation}\label{eq:gap-principle}
 s>\tfrac12q^{1+\delta}.
\end{equation}
For $\delta=a/b>0$ and $q\geq4^b$, this implies $s>2q$.
The identical argument can be conducted with rational positive certificates
and integer powers.  The gap principle forbids clustering and makes the
first successive denominators grow rapidly.  It does not exclude an
arbitrarily sparse infinite sequence; a global finite bound requires the
additional Roth estimates.

There is nevertheless useful constructive content in a certified
cardinality bound.  Suppose a decidable predicate $G(n)$ has at most $N$
solutions, and let $F$ be a supplied primitive recursive function with
$F(t)\geq t$.  Define
\[
 B_0=1,\qquad B_{i+1}=F(B_i)+1\quad(0\leq i\leq N).
\]
The $N+1$ disjoint finite intervals $[B_i,B_{i+1})$ cannot each contain a
solution.  Bounded searches find an empty one.  Thus there is an explicit
primitive recursive bound, depending on the supplied $F$, on a location of
a finite window with no solutions.  This is a valid route from a
quantitative Roth theorem to a metastable formulation, provided the
cardinality theorem and its arithmetic proof are actually supplied.
It still supplies no final cutoff valid for all later denominators.

\subsection{Lean and object-theory proof obligations}

The proved fragment can be organized into the following modules.
\begin{enumerate}
\item \texttt{RationalArithmetic}: signed integer and positive-denominator
codes; addition, multiplication, absolute value, powers, and decidable order.
Expose the cleared-denominator correctness lemmas.
\item \texttt{PolynomialDifference}: prove the finite factorization of
$x^i-y^i$ by induction on $i$, then
\eqref{eq:polynomial-lipschitz} by induction over the coefficient list.
\item \texttt{NoRationalRoot}: verify \eqref{eq:no-root-certificate}; prove
that $N_f(p,q)\ne0$ for arbitrary $p,q$ by rational reduction and the
rational-root argument.  The Euclidean algorithm needs a declared bounded
implementation or an explicit primitive recursive majorant.
\item \texttt{RootBisection}: define the bounded bisection loop on rational
endpoints and prove nesting, sign preservation, width, and residual bounds.
\item \texttt{LiouvilleFinite}: prove \eqref{eq:finite-separation} directly from the integer gap,
the difference bound, and the bisection residual bound at the same precision.
\item \texttt{QuadraticRoth}: define \eqref{eq:good-certificate} and prove
\eqref{eq:cutoff-explicit} by integer power arithmetic.
\end{enumerate}

The object-theoretic statements are schemes for fixed primitive recursive
operations, or use one fixed primitive recursive evaluator for the stated
finite polynomial and bisection data.  They do not posit a primitive
recursive evaluator for arbitrary primitive recursive programs.  Bounded
quantifiers in the input checks are implemented by Boolean folds; induction
in the argument is on explicit finite loops and can be presented as
quantifier-free induction on their invariants after coding finite records.
One must still translate the resulting derivations into the repository's
chosen iPRA calculus.  A Lean theorem over \texttt{Real} or a successful
\texttt{ring}/\texttt{norm\_num} invocation, by itself, would not accomplish
that translation.

For the full Roth target, a formalization plan must list the auxiliary
polynomial construction, coefficient estimates, index bounds, quantitative
selection of approximants, and extraction of the selected constructive
conclusion as separate proof obligations.  Calling any one of these an
unproved axiom would produce a conditional development.  A complete Lean
proof of the cutoff version would first require solving the mathematical
effectivity problem; the present report claims neither that result nor its
formalization.

\section{St\o rmer's theorem: a complete primitive recursive proof}
\label{sec:stormer}

This section proves the full finiteness-and-enumeration conclusion for consecutive
smooth integers.  The original St\o rmer method, described in the introduction to
Lehmer's paper~\cite{Lehmer1964}, uses coefficients whose prime exponents belong to
$\{0,1,2\}$.  Keeping the exponent $2$, rather than insisting on squarefree
coefficients, is decisive: it makes every prime in the Pell solution's second
coordinate divide the coefficient.  We give the elementary argument and a very
coarse explicit bound on the requisite fundamental solutions.  The bound is chosen
for a transparent proof in finite arithmetic, not for efficient numerical work.

\subsection{Statement and conventions}

Let $S=\{p_1,\ldots,p_k\}$ be a finite set of distinct primes.  A positive integer
is \emph{$S$-smooth} if all its prime factors belong to $S$; thus $1$ is smooth.
By convention $0$ is not smooth.
Put $R=\prod_{p\in S}p$, with the empty product equal to $1$, and define
\begin{equation}\label{eq:stormer-B}
 B(D)=(2D+2)^{\,2(2D+1)^3+2}\qquad(D\ge1).
\end{equation}
All exponents in this formula are natural-number exponents.

\begin{theorem}[Primitive recursive St\o rmer theorem]
\label{thm:stormer-pr}
There is a primitive recursive procedure which, from a finite list of distinct
primes $S$, returns precisely all positive integers $n$ for which $n$ and $n+1$
are both $S$-smooth.  Its output has at most $3^k-2^k$ entries.  Every output
satisfies
\[
             2n+1<B(R^2).
\]
If $2\notin S$, the output is empty.
\end{theorem}

We will explicitly prove every ingredient needed for this statement.  Square-root
notation in the proof is shorthand for exact arithmetic with pairs
$a+b\sqrt D$.  For nonsquare natural $D$, the sign of such a pair with rational
coefficients is decided by the signs of $a,b$ and a comparison of $a^2$ with
$Db^2$.  Alternatively, all occurrences of $\lfloor\sqrt m\rfloor$ are the
bounded-search function
\[
 \operatorname{isqrt}(m)=\max\{r\le m:r^2\le m\}.
\]
There is no real-number oracle or appeal to completeness.

\subsection{A bounded Pell construction}

\begin{lemma}[Finite approximation of a quadratic radical]
\label{lem:stormer-dirichlet}
For nonsquare $D\ge2$ and natural $N\ge1$, bounded search constructs positive
integers $p,q$, with $1\le q\le N$, such that
\[
                 |p-q\sqrt D|<1/N.
\]
\end{lemma}

\begin{proof}
For $0\le h\le N$ put
\[
 a_h=\operatorname{isqrt}(Dh^2),\qquad
 b_h=\operatorname{isqrt}(DN^2h^2)-Na_h.
\]
These integers satisfy
\[
 a_h\le h\sqrt D<a_h+1,
 \qquad
 b_h\le N(h\sqrt D-a_h)<b_h+1,
\]
so $0\le b_h<N$.  Among $N+1$ entries taking $N$ possible values, bounded
search finds $i<j\le N$ with $b_i=b_j$.  Set $q=j-i$ and $p=a_j-a_i$.
The difference of the two fractional parts has absolute value less than $1/N$,
which gives the required approximation.  Since $\sqrt D>1$ and $q\ge1$, the
integer $p$ is positive.  Finite pigeonhole here is the elementary counting fact
for bounded lists, proved by induction on the number of boxes.
\end{proof}

\begin{lemma}[An explicit bounded Pell solution]
\label{lem:stormer-pell-bound}
If $D\ge2$ is nonsquare, then bounded search finds positive integers $x,y$ such
that
\[
 x^2-Dy^2=1,
 \qquad x<B(D),\quad y<B(D).
\]
\end{lemma}

\begin{proof}
Write
\[
 C=2D+1,\quad L=C+1,\quad K=2C^3,\quad N_j=L^{j+1}
                    \quad(0\le j\le K).
\]
Apply Lemma~\ref{lem:stormer-dirichlet} with $N=N_j$ and denote its output by
$(p_j,q_j)$.  Put
\[
 t_j=p_j^2-Dq_j^2,\qquad U_j=p_j+q_j\sqrt D.
\]
The integer $t_j$ is nonzero: otherwise $D$ would be a square, by unique
factorization applied to $p_j^2=Dq_j^2$.  Moreover
\[
 U_j<2Dq_j+1\le Cq_j,
 \qquad
 0<|t_j|=|p_j-q_j\sqrt D|\,U_j<C.
\]
The displayed strict upper estimate for $U_j$ follows from $\sqrt D<D$ and
$|p_j-q_j\sqrt D|<1/N_j\le1$.
Because $|t_j|\ge1$, the same estimates give
\[
 \frac1{Cq_j}<\frac1{N_j},\qquad q_j>N_j/C.
\]
Consequently $q_{j+1}>N_{j+1}/C>N_j\ge q_j$.

Color the $K+1$ pairs by
\[
       \bigl(t_j,\ p_j\bmod |t_j|,\ q_j\bmod |t_j|\bigr).
\]
There are at most
$2\sum_{t=1}^{C-1}t^2\le2C^3=K$ colors.  A second finite pigeonhole search
therefore gives $i<j$ with the same color.  Write their common nonzero norm as
$t$.  Define
\[
 x=\frac{p_jp_i-Dq_jq_i}{t},\qquad
 y=\frac{q_jp_i-p_jq_i}{t}.
\]
Both quotients are integers: the first numerator is congruent modulo $|t|$ to
$p_i^2-Dq_i^2=t$, and the second is congruent to zero.  Multiplication in the
quadratic ring gives
\[
 x+y\sqrt D=\frac{U_j}{U_i},\qquad x^2-Dy^2=1.
\]
Since $q_j>q_i$ and the norms agree, $p_j>p_i$, so $U_j>U_i>1$.
For a norm-one number $V=U_j/U_i>1$,
\[
 x=\frac{V+V^{-1}}2>1,\qquad
 y=\frac{V-V^{-1}}{2\sqrt D}>0.
\]
Thus $x,y$ are positive integers, and
\[
 x<V<U_j<CN_j\le C L^{K+1}<L^{K+2}=B(D),
 \qquad y<V<B(D).
\]
All searches just specified have explicit natural-number bounds.
\end{proof}

\subsection{The structure of the positive Pell solutions}

Let $(u,v)$ be the positive solution of $u^2-Dv^2=1$ with least $u$; it can be
found by searching $1\le y<B(D)$, setting $x=\operatorname{isqrt}(Dy^2+1)$,
and taking the first exact solution.  Minimizing $y$ is equivalent to minimizing
$x$.  Set $\varepsilon=u+v\sqrt D$ and define integer pairs recursively by
\begin{align*}
 x_0&=1,&y_0&=0,\\
 x_{m+1}&=ux_m+Dvy_m,&y_{m+1}&=vx_m+uy_m.
\end{align*}
Then $x_m+y_m\sqrt D=\varepsilon^m$, and each pair has norm one.

\begin{lemma}[Bounded generation by the fundamental solution]
\label{lem:stormer-pell-powers}
Every positive solution $(x,y)$ of $x^2-Dy^2=1$ equals $(x_m,y_m)$ for a unique
positive index $m\le x$.
\end{lemma}

\begin{proof}
The coordinates of positive norm-one solutions are ordered together:
$x'<x$ if and only if $y'<y$, and then
$x'+y'\sqrt D<x+y\sqrt D$.  Since $u\ge2$, the recursion shows that
$x_m\ge m+1$ and that $x_m$ is strictly increasing.  Bounded search therefore
finds the largest $m\le x$ such that $x_m\le x$; necessarily $m<x$, and
\[
     \varepsilon^m\le x+y\sqrt D<\varepsilon^{m+1}.
\]
The quotient
\[
 \gamma=(x+y\sqrt D)(x_m-y_m\sqrt D)
\]
has integer coordinates, norm one, and $1\le\gamma<\varepsilon$.
If $\gamma>1$, both its norm-one coordinates are positive, and its first
coordinate is smaller than $u$, contrary to the definition of $(u,v)$.
Indeed, a first coordinate at least $u$ would force its positive second
coordinate to be at least $v$ by the norm equation, and hence force
$\gamma\geq\varepsilon$.
The decidable alternative is $\gamma=1$, yielding $(x,y)=(x_m,y_m)$.
The strict increase of $x_m$ proves uniqueness.  The index cannot be zero
because $y>0$.
\end{proof}

\subsection{St\o rmer's key divisibility lemma}

\begin{lemma}[Smooth second coordinate forces the fundamental solution]
\label{lem:stormer-key}
Suppose $x^2-Dy^2=1$ with $x,y>0$ and every prime divisor of $y$ divides $D$.
Then $(x,y)$ is the fundamental positive Pell solution for $D$.
\end{lemma}

\begin{proof}
By Lemma~\ref{lem:stormer-pell-powers}, write $(x,y)=(x_m,y_m)$.
We show $m=1$.  Suppose $m>1$ and choose a prime $\ell\mid m$, by bounded
trial division.  The binomial expansion of
$(x_\ell+y_\ell\sqrt D)^{m/\ell}$ shows that $y_\ell\mid y_m$.
Another binomial expansion gives $y_\ell=vC_\ell$, where
\begin{equation}\label{eq:stormer-binomial}
 C_\ell=\sum_{j=0}^{\lfloor(\ell-1)/2\rfloor}
  {\ell\choose 2j+1}u^{\ell-2j-1}v^{2j}D^j.
\end{equation}
Thus every prime dividing $C_\ell$ divides $D$.
The Pell identity gives $\gcd(u,D)=1$.
For any prime $r\mid C_\ell$, reduction of
\eqref{eq:stormer-binomial} modulo $r\mid D$ gives
\[
                   0\equiv\ell u^{\ell-1}\pmod r.
\]
Therefore $r=\ell$.  Since $C_\ell>1$, unique factorization yields
$C_\ell=\ell^e$ for some $e\ge1$ and also $\ell\mid D$.
We now examine the three decidable cases.

If $\ell=2$, then $C_2=2u$.  As $u>1$, a prime divisor of $u$ divides $C_2$
and hence $D$, contradicting $\gcd(u,D)=1$.

If $\ell=3$, the Pell identity reduces the expression to
\[
 C_3=3u^2+Dv^2=4u^2-1=(2u-1)(2u+1).
\]
Both factors exceed one, both are odd, and their greatest common divisor is one
because their difference is two.  They therefore cannot both be powers of
three, contradicting $C_3=3^e$.

Finally let $\ell\ge5$.  In \eqref{eq:stormer-binomial}, every term after the
first is divisible by $\ell^2$.  Indeed, the $j=1$ term contains both $D$ and
${\ell\choose3}$, each divisible by $\ell$; every term with $j\ge2$ contains
$D^2$.  Thus
\[
              C_\ell\equiv\ell u^{\ell-1}\pmod{\ell^2}.
\]
Since $\ell\mid D$ and $\gcd(u,D)=1$, the right-hand side is not divisible
by $\ell^2$.  Hence $e=1$ and $C_\ell=\ell$.
But the first, positive summand of \eqref{eq:stormer-binomial} is
$\ell u^{\ell-1}>\ell$, a contradiction.  Each possible prime divisor of
$m>1$ has been excluded, so $m=1$.
\end{proof}

Notice the special treatment of $\ell=3$.  A claim that the same valuation-one
argument works for every odd prime would be false: the binomial coefficient
${3\choose3}$ contributes no extra factor of three.

\subsection{Reduction of smooth neighbors and the finite algorithm}

\begin{proof}[Proof of Theorem~\ref{thm:stormer-pr}]
If $2\notin S$, every $S$-smooth number is odd, so there are no consecutive
pairs.  Assume henceforth that $2\in S$ and that $n,n+1$ are $S$-smooth.
Factor
\[
             (2n+1)^2-1=4n(n+1)=\prod_{i=1}^k p_i^{a_i}.
\]
Define
\[
 e_i=\begin{cases}
 0,&a_i=0,\\
 1,&a_i\text{ is odd},\\
 2,&a_i>0\text{ is even},
 \end{cases}
 \qquad
 D=\prod_i p_i^{e_i},\qquad
 y=\prod_i p_i^{(a_i-e_i)/2}.
\]
Each exponent is a nonnegative integer.  We have
\[
 (2n+1)^2-Dy^2=1,
\]
and every prime dividing $y$ divides $D$, since positive $a_i$ always gives
positive $e_i$.  The coefficient $D$ is nonsquare: if $D=d^2$, then
\[
 (2n+1-dy)(2n+1+dy)=1
\]
would force $dy=0$, impossible.  Lemma~\ref{lem:stormer-key} now says that
$(2n+1,y)$ is the fundamental solution for $D$.

There are exactly $3^k$ possible coefficient vectors $(e_1,\ldots,e_k)$.
Exactly $2^k$ give squares, namely those with every $e_i\in\{0,2\}$.
A nonsquare coefficient has only one fundamental solution and therefore yields
at most one neighbor pair.  This proves the bound $3^k-2^k$.
Also $D\le R^2$ and $B$ is increasing, so
\[
                 2n+1<B(D)\le B(R^2).
\]

Here is a complete algorithm.  Enumerate all vectors in $\{0,1,2\}^k$ and
form their coefficients $D$.  Discard squares.  For each remaining $D$, search
$1\le y<B(D)$ until $x=\operatorname{isqrt}(Dy^2+1)$ satisfies
$x^2-Dy^2=1$.  The preceding bounded Pell lemma guarantees success.
If $x$ is odd, set $n=(x-1)/2$ and retain it precisely when both $n$ and $n+1$
are positive and $S$-smooth.  Remove duplicate entries and sort.
The construction just proved shows that every required pair occurs;
the final tests show that every returned pair is correct.  The loops are bounded
by explicitly primitive recursive functions.  No stopping test for an infinite
search is being assumed.
\end{proof}

For $S=\{2,3,5\}$ this returns
\[
 (1,2),\ (2,3),\ (3,4),\ (4,5),\ (5,6),\ (8,9),\ (9,10),
 \ (15,16),\ (24,25),\ (80,81).
\]
The accompanying program can compute such small examples while retaining the
proved bound as a totality bound.  The large value of $B(D)$ affects the worst
case bound, not the early return when the fundamental solution is small.

\subsection{Why the proof is available in iPRA}

The mathematical content above reduces to a uniform assertion about finite
codes.  Trial division, primality testing, integer square root, signed integer
arithmetic, exponentiation, binomial coefficients, finite list coding, and
bounded minimization are primitive recursive.  A $D$-smoothness assertion is a
bounded divisibility test.  Algebraic comparisons in $\mathbb Q(\sqrt D)$ are
decidable by integer squaring and sign cases.  The two pigeonhole applications
are finite counting identities.  The power-sequence identities follow by
quantifier-free induction on the index; binomial divisibility follows by the
usual integer product identities and Euclid's lemma.  The apparent universal
properties of lists are bounded predicates and can be replaced by their
primitive recursive Boolean characteristic functions.

For complete proof-theoretic precision, the iPRA assertion can be stated as a
schema for fixed choices of these primitive recursive symbols, with free
natural-number inputs coding $S$ and $n$:
\[
 \operatorname{PrimeList}(S)\ \Longrightarrow\
 \Bigl[\bigl(\operatorname{Smooth}(S,n)\wedge
       \operatorname{Smooth}(S,n+1)\bigr)
 \ \Longleftrightarrow\ n\in\operatorname{Neighbors}(S)\Bigr].
\]
The membership predicate on the right is bounded in the explicitly constructed
output list.  The auxiliary cardinality and height statements are likewise
finite arithmetic statements.  A fully verified implementation must still
supply the iPRA derivations of the lemmas; an ordinary Lean proof over natural
numbers alone is not automatically such a derivation.

\subsection{Lean implementation path}

An implementation can be divided into the following independently checkable
parts.  First encode ordered quadratic pairs and prove the integer sign test,
ring identities, norm identities, and the formulas used by the quotient
construction.  Next formalize the two bounded pigeonhole searches and prove the
output specification of the Pell constructor.  Then prove generation by the
fundamental solution, the displayed binomial formula, and the three prime-index
cases.  Finally define the coefficient enumeration and smoothness filters,
and prove soundness, completeness, the cardinal bound, and the explicit height
bound.  Every unbounded-looking existential in this route has already been
replaced above by a stated finite search bound.

For an iPRA certificate rather than only a metatheoretic Lean proof, use the
repository's primitive recursive syntax and derivation infrastructure to
represent those definitions, and translate the displayed bounded induction
invariants into characteristic-function equations.  The substantive remaining
formalization work is the checked translation and arithmetic infrastructure;
there is no omitted Diophantine estimate or primitive-divisor theorem in this
St\o rmer argument.

\section{A constructive Davenport--Schmidt theorem with an explicit modulus}
\label{sec:ds}

The theorem linked in the question concerns approximation by algebraic
numbers of degree at most two, with exponent three.  It is not the later
Davenport--Schmidt theorem about simultaneous approximation to a number and
its square.  The original 1967 result gives the constant
$C>160/9$ in
\[
  |\xi-\alpha|< C\max(1,\xi^2)H(\alpha)^{-3}.
\]
Here we retain the exponent and the dependence on $\xi$, but use the
larger absolute constant $32768$.  In exchange, every search below has an
explicit primitive recursive bound.  The proof is a finite reconstruction
of the minimal-point argument of Davenport and Schmidt
\cite{DavenportSchmidt1967}, with an elementary pigeonhole estimate in
place of its sharper geometry-of-numbers estimate.

\subsection{The positive hypothesis and the exact claim}

For a nonzero integer polynomial $f=aX^2+bX+c$ of degree at most two put
$\|f\|=\max(|a|,|b|,|c|)$.  A quadratic separation modulus for $\xi$ is a
nondecreasing function $\rho:\mathbb N_{\geq1}\to\mathbb N$ such that
\begin{equation}\label{eq:ds-separation}
  0<\|f\|\leq H\quad\Longrightarrow\quad
  |f(\xi)|\geq 2^{-\rho(H)}.
\end{equation}
This is positive information.  It implies that $\xi$ is neither rational
nor quadratic.  We do not infer the existence of a primitive recursive
$\rho$ from the negative assertion that no quadratic polynomial vanishes.

\begin{theorem}[Constructive Davenport--Schmidt, with supplied separation]
\label{thm:ds-constructive}
Fix a primitive recursive rational name for a real number $\xi$, a
primitive recursive modulus $\rho$ satisfying
\eqref{eq:ds-separation}, and an integer $M\geq1$ with $|\xi|\leq M$.
For every integer $h\geq1$, a bounded primitive recursive procedure
produces a rational number or real quadratic irrational $\alpha$ such that
\[
 H(\alpha)>h,\qquad
 |\xi-\alpha|<32768\max(1,\xi^2)H(\alpha)^{-3}.
\]
The output consists of its primitive minimal polynomial and a root
selector.  All bounds used by the procedure are displayed below.

In first-order iPRA this is a schema for the fixed primitive recursive
name and modulus, conditional on their stated arithmetical correctness
properties.  It does not quantify over arbitrary functions as objects.
\end{theorem}

An algebraic real of fixed degree at least three admits such a modulus by
an elementary resultant estimate.  Thus the hypothesis has substantial
unconditional instances; it is not restricted to a hypothetical oracle.
A construction of the modulus is given at the end of this section.

We first prove a finite arithmetic core at a rational working point
$0<\eta<1$.  Temporarily, $\rho$ denotes the separation exponent used
by this core.  Separation is required only for coefficient heights at most
$W=2B^2$, where $B$ is the finite bound constructed below.  This is
consistent with $\eta$ being rational: its own linear relation may have
height much greater than $W$.  In the actual algorithm we will compute
$B$ with the exponent $\rho+1$ and obtain the working point by one
sufficiently accurate rational sample of the given real number.  Write
\[
 K=16,\qquad D=4096,\qquad C=8192=2D,\qquad
 \varepsilon(H)=2^{-\rho(H)}.
\]
The rational core has constant $C$.  Its rational sample will give the
normalized real conclusion with constant $16384$.

\subsection{A finite pigeonhole estimate}

For every integer $3\leq U\leq B$ there is a nonzero integer triple
$v=(a,b,c)$ with
\begin{equation}\label{eq:ds-dirichlet}
 \|v\|<U,\qquad 0<|a\eta^2+b\eta+c|<\frac{16}{U^2}.
\end{equation}
To prove this, put $N=\lfloor(U-1)/2\rfloor$.  The $(N+1)^2$ numbers
\[
 a\eta^2+b\eta\qquad(0\leq a,b\leq N)
\]
have fractional parts in $[0,1)$.  Partition that interval into $N^2$
half-open intervals of length $N^{-2}$.  Two distinct pairs give fractional
parts in the same interval.  Subtract them, and subtract the corresponding
integer parts.  The resulting $a,b$ have absolute values at most $N$;
the constant coefficient has absolute value at most $2N$ because the
original numbers lie in $[0,2N]$.  The triple is nonzero, since the pairs
are distinct.  Its value is nonzero by the separation hypothesis and has
absolute value less than $N^{-2}$.  Finally $2N<U$ and $N\geq U/4$.

This pigeonhole argument is constructive.  At the rational working point,
integer parts and occupied subintervals are computed by exact integer
division.  Boundary equality is decidable.  The ordinary finite
pigeonhole proof uses induction only on finite lists.  The separation
assumption is needed only for the final nonzero triple, whose height is
less than $U\leq B$.

\subsection{Finite minimal points and a bound for the next record}

For each integer $1\leq H\leq B$, consider the finite collection of nonzero
integer triples of height at most $H$.  After changing signs when needed,
their values under
\[
 L(a,b,c)=a\eta^2+b\eta+c
\]
are positive.  Choose the triple having the smallest positive value.
It is unique: equality of two values would give a nonzero quadratic
polynomial of height at most $2B\leq W$ vanishing at $\eta$.  List the distinct choices in increasing
order of the height at which they first occur, stopping at any specified
finite bound.  Denote them by $v_i$, and put
\[
 H_i=\|v_i\|,\qquad E_i=L(v_i)>0,\qquad
 P(v)=2\eta a+b,\qquad P_i=|P(v_i)|.
\]
Then $H_i$ strictly increases, $E_i$ strictly decreases, and every $v_i$
is primitive.  Indeed, a common divisor greater than one would give a
smaller positive value at a smaller height.  Consecutive $v_i$ are linearly
independent: primitive proportional integer vectors agree up to sign,
which is incompatible with their distinct positive values.

Whenever $H_{i+1}\geq3$, \eqref{eq:ds-dirichlet} applied with
$U=H_{i+1}$ gives
\begin{equation}\label{eq:ds-record-estimate}
 E_i<\frac{K}{H_{i+1}^2}.
\end{equation}
The reason is that $E_i$ is the least positive absolute value at heights
strictly below $H_{i+1}$.

Define the primitive recursive function
\begin{equation}\label{eq:ds-jump}
 J(R)=\max\{R+2,\,8\cdot2^{\rho(R)}\}\qquad(R\geq1).
\end{equation}
Whenever $J(R)\leq B$, there is a record of height strictly greater
than $R$ and strictly less than $J(R)$.  In fact, \eqref{eq:ds-dirichlet} at $J(R)$ supplies a value
less than
\[
 \frac{16}{J(R)^2}\leq\frac14\,2^{-2\rho(R)}
 <2^{-\rho(R)},
\]
whereas every nonzero triple of height at most $R$ has value at least
$2^{-\rho(R)}$.  This proves a bound for obtaining a later record without
any appeal to an infinite sequence or an unbounded search.

\subsection{The integral-plane lemma}

Two consecutive minimal points $v_i,v_{i+1}$ form an integral basis for
the integer points in their rational plane.  To see this, if an integer
point in that plane were outside their integral span, reduce its two
rational coordinates modulo integers into $[-1/2,1/2]$.  This produces
a nonzero integer point
\[
 w=r v_i+s v_{i+1},\qquad |r|,|s|\leq\tfrac12.
\]
It satisfies
\[
 \|w\|\leq\tfrac12(H_i+H_{i+1})<H_{i+1},\qquad
 |L(w)|\leq\tfrac12(E_i+E_{i+1})<E_i,
\]
contradicting minimality.  The coordinates can be computed by a nonzero
two-by-two minor, so this argument involves only rational arithmetic.

If three consecutive points are linearly dependent, the two adjacent
pairs are integral bases of the same plane.  Consequently
\begin{equation}\label{eq:ds-plane-relation}
 v_{i-1}=u v_i+s v_{i+1},\qquad u\in\mathbb Z,\quad s\in\{1,-1\}.
\end{equation}
For completeness, expansion in one integral basis gives the coefficient
$s$ as an integer.  Expansion in the other basis shows that its inverse
is an integer.  Hence $s=\pm1$.
For any two linear forms $A,B$, relation
\eqref{eq:ds-plane-relation} preserves, up to sign, the quantity
\[
 A(v_{i-1})B(v_i)-A(v_i)B(v_{i-1}).
\]
In particular, its absolute value is constant along a finite string of
dependent consecutive triples.

\subsection{A finite bound for finding an independent triple}

Suppose a pair of consecutive records $v_r,v_{r+1}$ has heights at most
$R$, and the upper bound $J(G(R))$ below is at most $B$.  Some coordinate $j$ makes
\[
 w=(v_r)_j v_{r+1}-(v_{r+1})_j v_r
\]
nonzero, because the pair is independent.  We have $\|w\|\leq2R^2$, and
therefore
\begin{equation}\label{eq:ds-wedge-lower}
 |L(w)|\geq2^{-\rho(2R^2)}.
\end{equation}
Define
\begin{equation}\label{eq:ds-gap}
 G(R)=\max\{R+1,\,32\cdot2^{\rho(2R^2)}+1\}.
\end{equation}
Take a later record $v_t$ with $H_t>G(R)$, available below $J(G(R))$.
If every triple with center $r+1,r+2,\ldots,t-1$ were dependent, the
integral-plane lemma would make the absolute value in
\eqref{eq:ds-wedge-lower} equal to
\[
 |(v_{t-1})_jE_t-(v_t)_jE_{t-1}|
 \leq H_{t-1}E_t+H_tE_{t-1}
 <\frac{2K}{H_t}
 <2^{-\rho(2R^2)}.
\]
The penultimate inequality uses \eqref{eq:ds-record-estimate} and
$E_t<E_{t-1}$.  This contradiction proves that an independent triple
occurs in the displayed finite range.  No assertion about infinitely many
independent triples has been used.

\subsection{The determinant calculation}

We now prove the key finite obstruction.  Suppose, over the finite range
under consideration, that every relevant record satisfies
\begin{equation}\label{eq:ds-small-derivative}
 P_i\leq\frac{H_i^3E_i}{D}.
\end{equation}
Let $n$ be the center of an independent triple.  The determinant with
rows
\[
 ((v_i)_1,P(v_i),L(v_i)),\qquad i=n-1,n,n+1,
\]
equals the determinant with rows $v_i$: the triangular transformation
from $(a,b,c)$ to $(a,2\eta a+b,\eta^2a+\eta b+c)$ has determinant one.
It is thus a nonzero integer and has absolute value at least one.

Expand along the middle column.  The term involving $P_{n-1}$ has absolute
value at most
\[
 P_{n-1}(H_nE_{n+1}+H_{n+1}E_n)
 <\frac{2}{D}H_{n+1}H_{n-1}^3E_{n-1}E_n
 <\frac{2K^2}{D}=\frac18.
\]
Here the last inequality follows from
$E_{n-1}<K/H_n^2$, $E_n<K/H_{n+1}^2$, and
$H_{n-1}<H_n<H_{n+1}$.  The term involving $P_n$ is similarly less than
\[
 \frac{2}{D}H_{n+1}H_n^3E_{n-1}E_n
 <\frac{2K^2}{D}=\frac18.
\]
The remaining cofactor has absolute value at most $2H_nE_{n-1}$.
Therefore
\begin{equation}\label{eq:ds-derivative-lower}
 P_{n+1}H_nE_{n-1}>\frac38,
 \qquad P_{n+1}>\frac{3H_n}{128}.
\end{equation}
The second inequality follows again from $E_{n-1}<16/H_n^2$.

Let $m<n$ be consecutive centers of independent triples.  The intervening
triples are dependent, so the integral-plane lemma gives
\begin{align*}
 \Omega&=|P(v_m)E_{m+1}-P(v_{m+1})E_m|\\
  &=|P(v_{n-1})E_n-P(v_n)E_{n-1}|.
\end{align*}
The latter expression is bounded above by
\[
 \frac1D(H_{n-1}^3+H_n^3)E_{n-1}E_n
 <\frac{2K}{D}H_nE_n=\frac{2}{256}H_nE_n.
\]
The former, by \eqref{eq:ds-derivative-lower}, is bounded below by
\[
 \frac{3}{128}H_mE_m-\frac1D H_m^3E_mE_{m+1}
 >\left(\frac{3}{128}-\frac{K}{D}\right)H_mE_m
 =\frac{5}{256}H_mE_m.
\]
We used $E_{m+1}<E_m<K/H_m^2$, so this estimate does not need a further
successor beyond $v_{m+1}$.  Combining the two bounds gives
\begin{equation}\label{eq:ds-growth}
 H_nE_n>\frac52 H_mE_m.
\end{equation}
This is a finite growth invariant, stronger than the monotonicity needed
below.

\subsection{An explicit height bound for a large derivative}

Let $T\geq16385$.  Define, in the following order,
\begin{equation}\label{eq:ds-final-bound}
\begin{aligned}
 A&=J(T),& A'&=J(A),& U&=J(G(A')),\\
 R&=\max\{U+1,16\cdot2^{\rho(U)}+1\},\\
 V&=J(R),& V'&=J(V),& B&=J(G(V')).
\end{aligned}
\end{equation}
All of these are primitive recursive functions of $T$ for the fixed
primitive recursive function $\rho$.

\begin{lemma}\label{lem:ds-large-derivative}
Among the primitive nonzero integer triples with $T<\|v\|\leq B$ there is
one satisfying
\begin{equation}\label{eq:ds-large-derivative}
 \|v\|^3|L(v)|<D|P(v)|.
\end{equation}
\end{lemma}

\begin{proof}
Assume the finite, decidable search for such a triple returns no match.
Then \eqref{eq:ds-small-derivative} holds for every record of height
strictly between $T$ and $B$.

The first record above $T$ has height below $A$; its successor has height
below $A'$.  The finite independence lemma supplies an independent triple
whose center $m$ is after that pair's first point and has $H_m<U$.
In particular, all three points involved have height greater than $T$.
The separation hypothesis and monotonicity of $\rho$ give
\[
 H_mE_m\geq E_m\geq2^{-\rho(U)}.
\]
Likewise the first record above $R$ and its successor have heights below
$V$ and $V'$.  Applying the finite independence lemma again supplies an
independent center $n$ with
\[
 R<H_n<B,
\]
whose successor is also within the searched range.  Apply
\eqref{eq:ds-growth} successively to the finite list of independent
centers from $m$ to $n$.  It follows that
$H_nE_n>2^{-\rho(U)}$.  But
\[
 H_nE_n<\frac{16}{H_n}
 <2^{-\rho(U)}
\]
by \eqref{eq:ds-record-estimate} and the definition of $R$.  This is a
contradiction.

The search is over a finite decidable list.  Its failure leads to a
contradiction, so its other, explicitly recorded outcome supplies a
witness.  This last step is finite decidable case analysis, not an
application of double-negation elimination for an unbounded existential.
\end{proof}

\subsection{Turning the large derivative into a nearby algebraic number}

Let $f=aX^2+bX+c$ be the primitive polynomial supplied by
Lemma~\ref{lem:ds-large-derivative}, and write $H=\|f\|$,
$L=f(\eta)$, $P=f'(\eta)$.  Then $L\ne0$, $P\ne0$, $H>4D$, and
\[
 |L|<D|P|H^{-3}.
\]
We claim that $f$ has a real root $\alpha$ with
\begin{equation}\label{eq:ds-root-distance}
 |\eta-\alpha|<2DH^{-3}.
\end{equation}

If $|P|\geq H^{-1}$, put $t=2DH^{-3}$ and choose
$s=-\operatorname{sgn}(LP)t$.  The exact quadratic Taylor identity gives
\[
 \operatorname{sgn}(L)f(\eta+s)
 \leq |L|-|P|t+H t^2
 <-D|P|H^{-3}+4D^2H^{-5}<0.
\]
The last inequality uses $H>4D$.  Thus the values of $f$ at $\eta$ and
$\eta+s$ have opposite signs, and its explicit linear or quadratic root
formula supplies a root strictly between them.  This use of a sign
change needs no general intermediate-value principle: it is just the
elementary quadratic formula, or equivalently its elementary algebraic
sign verification.  Indeed, throughout $|z|\leq t$,
\[
 |f'(\eta+z)-P|\leq2Ht=4D/H^2<1/H\leq|P|,
\]
so the derivative keeps the sign of $P$.  If $a\ne0$, the discriminant
$\Delta=b^2-4ac$ is positive, and the selected root is explicitly
\[
 \alpha=\frac{-b+\operatorname{sgn}(P)\sqrt{\Delta}}{2a}.
\]

If $|P|<H^{-1}$, the integer discriminant satisfies
\[
 |b^2-4ac|=|P^2-4aL|
 <H^{-2}+4DH^{-3}<1.
\]
Hence it is zero.  Here $a\ne0$, since otherwise $P=b=0$.
Primitivity implies
\[
 f(X)=\pm(uX+v)^2
\]
for coprime integers $u,v$ with $u\ne0$.  To see this, put the repeated
rational root in lowest terms $-v/u$.  Then
$f=a(X+v/u)^2$.  Since $c=av^2/u^2$ is an integer and
$\gcd(u^2,v^2)=1$, we have $a=ku^2$ for an integer $k$.  Thus the
coefficients are $ku^2,2kuv,kv^2$.  The three numbers
$u^2,2uv,v^2$ have greatest common divisor one, so primitivity forces
$k=\pm1$.  Since $u\eta+v\ne0$, the
large-derivative inequality reduces to
\[
 |u\eta+v|<2D|u|H^{-3}.
\]
The rational root $\alpha=-v/u$ again satisfies
\eqref{eq:ds-root-distance}.

In both cases $H(\alpha)\leq H$.  For an irreducible quadratic this is
equality.  For a rational root in lowest terms $p/q$, the rational-root
divisibility relations give $|p|\leq|c|$, $q\leq|a|$ when $ac\ne0$;
the cases of a zero constant coefficient and a linear polynomial are immediate.  Consequently
\[
 |\eta-\alpha|<8192H(\alpha)^{-3}.
\]

\subsection{Forcing the output height to exceed a prescribed height}

Fix $h\geq1$.  If $\beta$ has degree at most two and $H(\beta)\leq h$,
let $g$ be its primitive minimal polynomial.  For
$|\eta-\beta|<1$ we have $|\beta|<2$, and
\[
 |g(\eta)|
 =|\eta-\beta|\,|a(\eta+\beta)+b|
 \leq4h|\eta-\beta|.
\]
The same formula covers a linear $g$ by setting $a=0$.  Therefore
\begin{equation}\label{eq:ds-small-height-separation}
 |\eta-\beta|\geq\frac{2^{-\rho(h)}}{4h}.
\end{equation}
Set
\begin{equation}\label{eq:ds-threshold}
 T=\max\{16385,\,32768h\cdot2^{\rho(h)}+1\}.
\end{equation}
For every $H>T$, the root distance in
\eqref{eq:ds-root-distance} is less than both $1$ and the right hand side
of \eqref{eq:ds-small-height-separation}.  The produced root must therefore
have height greater than $h$.  Equations
\eqref{eq:ds-jump}, \eqref{eq:ds-gap}, \eqref{eq:ds-final-bound}, and
\eqref{eq:ds-threshold} are the promised explicit height-search bound.

\subsection{Freezing one rational sample: the strict iPRA version}

We now distinguish the given real $x\in(0,1)$ from the rational working
point $\eta$.  Let $\rho$ be the supplied separation exponent for $x$.
Put
\[
 \sigma(H)=\rho(H)+1.
\]
Compute $T,J,G,B$ from
\eqref{eq:ds-jump}, \eqref{eq:ds-gap}, \eqref{eq:ds-final-bound}, and
\eqref{eq:ds-threshold}, using $\sigma$ everywhere in those formulas.
Set $W=2B^2$.  Take one rational sample $\eta$ of the given name such that
\begin{equation}\label{eq:ds-freeze}
 |\eta-x|\leq
 \min\left\{\frac{2^{-\rho(W)}}{8W},\,B^{-3}\right\}.
\end{equation}
This is a primitive recursive precision request.  For example, a name
whose $s$th rational has error at most $2^{-s}$ can be sampled at
\[
 s=\rho(W)+8W+3B+10.
\]
The deliberately excessive additions avoid any need to compute logarithms.

The inequalities $x\geq2^{-\rho(1)}$ and
$1-x\geq2^{-\rho(1)}$ follow by applying separation to $X$ and $X-1$.
Consequently \eqref{eq:ds-freeze} ensures $0<\eta<1$.  For any polynomial
$f$ of height at most $H\leq W$, the elementary estimate
\[
 |f(\eta)-f(x)|\leq4H|\eta-x|
 \leq\tfrac12 2^{-\rho(W)}
 \leq\tfrac12 2^{-\rho(H)}
\]
shows
\begin{equation}\label{eq:ds-local-separation}
 |f(\eta)|\geq2^{-\sigma(H)}
 \quad(0<\|f\|\leq H\leq W).
\end{equation}
Thus the entire finite core can be run at this one rational $\eta$.
All of its separation uses lie inside this local range: differences of
searched triples have height at most $2B$; the integer vectors in the
independence lemma have height at most $2R^2$ with $R<B$; the next-record
and small-output-height arguments use heights below $B$.

The core gives a root $\alpha$ of a primitive polynomial of height
$H\leq B$, with $H(\alpha)>h$ and
$|\eta-\alpha|<8192H^{-3}$.  Therefore
\[
 |x-\alpha|
 \leq|x-\eta|+|\eta-\alpha|
 <B^{-3}+8192H^{-3}
 \leq8193H(\alpha)^{-3}
 <16384H(\alpha)^{-3}.
\]
The output-height exclusion was already proved for the rational working
point, so no separate limiting argument is needed.

This freezing step is also the foundation audit.  Every induction in
the core concerns finite integer or rational codes at the single rational
input $\eta$.  Values such as $P_iE_j-P_jE_i$ and $H_iE_i$ are rational
numbers, so their signs and inequalities are decidable by cross
multiplication; no hidden separation modulus for quartic expressions is
required.  Local condition \eqref{eq:ds-local-separation} is one bounded
primitive recursive predicate on the finite input.  Quantifier-free
induction is sufficient for the finite list invariants and bounded
searches.  Only the final error estimate refers back to the given name.

\subsection{Removing the normalization}

The separation modulus decides the signs of $\xi$, $\xi-1$, and
$\xi+1$; none is zero.  Reflection preserves algebraic heights.  Thus
if $0<|\xi|<1$, apply the normalized theorem to $x=|\xi|$ and then
reflect the output when necessary.

If $|\xi|>1$, put $x=1/|\xi|\in(0,1)$.  Reversal of the three
coefficients and multiplication by $\xi^2$ show that
\[
 |ax^2+bx+c|
 \geq M^{-2}2^{-\rho(H)}
 \geq2^{-(\rho(H)+2M)}.
\]
Thus $\rho_x(H)=\rho(H)+2M$ is an admissible primitive recursive
separation exponent for $x$.  A primitive recursive rational name
for $x$ is obtained by division, using $|\xi|>1$ as the denominator
bound.  Run the normalized construction with requested minimum height
\[
 h'=\max\{h,32768M\}.
\]
Its output $\beta$ satisfies
$|x-\beta|<16384H(\beta)^{-3}<1/(2M)$, so
$|\beta|>1/(2|\xi|)>0$.  Let $\alpha$ be the signed reciprocal of
$\beta$.  Reversal of a primitive minimal polynomial preserves its
height, and hence $H(\alpha)=H(\beta)>h$.  Finally
\[
 |\xi-\alpha|
 =\frac{|x-\beta|}{x|\beta|}
 <2\xi^2|x-\beta|
 <32768\xi^2H(\alpha)^{-3}.
\]
This completes the constructive proof of
Theorem~\ref{thm:ds-constructive}.

\subsection{Why the construction is primitive recursive}

Here are details that prevent the phrase ``effective'' from concealing
an unbounded search.  For the fixed primitive recursive $\rho$, all of
$T,J,G,B,W,s$ are explicit finite compositions of primitive recursive
functions.  After sampling once at precision $s$, enumerate the triples
in $[-B,B]^3$, test their content by Euclid's algorithm, and retain those
in the specified height range.  At this rational input every value,
minimum, comparison, record, and large-derivative test is exact rational
arithmetic.  No further refinement loop or sign oracle is used.

Quadratic root output requires only an integer square-root test and a
choice between two explicit roots.  The square root of a positive integer
has a primitive recursive rational name obtained by a prescribed number
of rational bisections.  Reduction of a rational root to lowest terms is
primitive recursive.  The root's primitive polynomial and sign selector
are finite data; approximation of its value to any requested precision
is a separate bounded primitive recursive calculation.

The iPRA proof of the rational core uses the following finite invariants:
the smallest positive value seen so far; the height and value ordering
of records; the integral-plane relation and its preserved two-by-two
minors; and the growth inequality \eqref{eq:ds-growth} on a finite list.
After multiplying by positive denominators, each is a predicate on
integer codes.  Bounded quantifiers are compiled into Boolean primitive
recursive tests before quantifier-free induction is applied.  The local
separation predicate is an explicit premise of this finite rational
theorem, and \eqref{eq:ds-freeze}--\eqref{eq:ds-local-separation} explain
how the fixed real name and modulus supply it.  The mere existence of a
terminating program would not, by itself, establish this translation;
the rational freezing is what makes its induction formulas explicit.

\subsection{An explicit separation modulus for algebraic inputs}

Suppose $\xi$ is a specified real root of an irreducible integer
polynomial $F$ of degree $d\geq3$, with coefficient height $A\geq1$ and
$|\xi|\leq R$, where $R\geq1$ is an integer.  For any nonzero integer
polynomial $g$ of degree $e\leq2$ and height at most $H$, its resultant
with $F$ is a nonzero integer.  The Sylvester-matrix adjugate supplies
integer polynomials $U,V$ with
\[
 UF+Vg=\operatorname{Res}(F,g),\qquad \deg V<d.
\]
Writing $M_H=\max(A,H,1)$, the Leibniz determinant formula bounds every
coefficient of $V$ by
\[
 (d+1)!M_H^{d+1}.
\]
Indeed the Sylvester matrix has size at most $d+2$, and its cofactors
are determinants of size at most $d+1$ with entries bounded by $M_H$.
For constant $g$, the direct lower bound $|g(\xi)|\geq1$ is even simpler.
Evaluation at $\xi$ yields
\[
 1\leq|\operatorname{Res}(F,g)|
 \leq|g(\xi)|\,d(d+1)!M_H^{d+1}(R+1)^{d-1}.
\]
Define
\[
 Q(H)=d(d+1)!\max(A,H,1)^{d+1}(R+1)^{d-1}.
\]
Then $|g(\xi)|\geq Q(H)^{-1}\geq2^{-Q(H)}$.
Thus $\rho(H)=Q(H)$ is an explicit nondecreasing primitive recursive
separation exponent.  The estimate is deliberately coarse.

For a fixed algebraic root one may use its usual primitive recursive
isolating-interval name.  To avoid concealing a termination assertion,
the input can instead include a rational interval whose endpoint
values of $F$ have opposite signs, together with a positive lower bound
on $|F'|$ throughout the interval; sufficiently small such intervals
exist for a simple algebraic root, and once supplied, bisection has an
explicit precision bound.  These finite certificates are checkable by
rational polynomial estimates.  Uniform construction of all algebraic
root data can be treated separately by a bounded root-isolation algorithm.

\subsection{Lean implementation plan and the remaining formal work}

The mathematical proof above is complete for the stated positive
hypotheses; no assertion is made that the entire iPRA derivation has
already been machine checked.  A suitable Lean development has these
modules, in dependency order:
\begin{enumerate}
 \item Integer triples, content, height, $L$, $P$, polynomial reversal,
       and finite list enumeration.  State the separation hypothesis for
       fixed encoded primitive recursive functions, and the bounded
       local separation predicate for a rational working point.
 \item One rational freezing operation with the explicit precision $s$,
       proving the bounded local separation predicate.  All later
       decisions use exact rational arithmetic.
 \item The finite pigeonhole lemma and the record algorithm, proving
       \eqref{eq:ds-record-estimate} and the next-record bound $J$.
 \item The integral-plane lemma using two-by-two minors and rational
       coordinate reduction; then the finite independence bound $G$.
 \item The determinant identity and the rational inequalities leading
       to \eqref{eq:ds-growth}.  These are fixed-degree polynomial
       identities and ordered rational arithmetic after denominators
       are cleared.
 \item The finite search bound $B$, proved by the two bounded
       applications of the independence lemma and finite induction over
       the list of independent centers.
 \item Linear/quadratic root certificates, the discriminant-zero
       factorization, minimal-polynomial height comparison, and the
       exclusion of heights at most $h$.
 \item Reflection and reciprocal transformations, and optionally the
       Sylvester-adjugate separation bound for algebraic inputs.
\end{enumerate}
Proving the corresponding assertions over Lean's real numbers is useful
for checking the mathematics, but it is not the final foundation audit.
The iPRA layer must prove the bounded algorithms' rational certificate
properties in the repository's syntax and exhibit the relevant
quantifier-free induction instances.  This is a clear, finite proof
engineering task; an appeal to classical real-number minimum selection
would lose the constructive content established here.

\section{Formalization milestones and computational evidence}
\label{sec:formalization}

\subsection{A route to actual iPRA derivations}

The supplied Lean files are deliberately explicit about their status.
\path{FiniteSpecifications.lean} gives concrete numerical bounds, rational
predicates, a finite Davenport--Schmidt target, and several elementary ring
identities.  \texttt{ObjectTargets.lean} gives the scoped first-order shape of
St\o rmer completeness and applies the repository's existing transport to an
already supplied derivation.  Neither file supplies the main missing
Diophantine derivations.  The environment had no Lean executable, so the files
have not been elaborated or kernel checked.  They target the inspected
repository toolchain, \texttt{leanprover/lean4:v4.32.0}.

The following milestones specify the work required to complete the
formalization.  They are ordered so that each has a concrete, independently
reviewable output.

\begin{enumerate}
\item \textbf{Arithmetic definitions and characteristic predicates.}
Choose the numerical codec.  Give actual \texttt{PR} syntax trees for signed
arithmetic, rational operations, finite lists, bounded search, powers, gcd,
integer square root, and polynomial evaluation.  Prove their defining
equations and arithmetic invariants in \texttt{Prf PRAAx}.  Bound the
Euclidean algorithm by its initial input and then present it as a clocked
primitive recursion.  Compile every bounded quantifier into a Boolean fold.

\item \textbf{Finite counting and finite selection.}
Formalize the occupancy-table proof of pigeonhole, with a returned pair of
indices.  Formalize bounded minima, records, and prime-divisor search.  Each
search returns both a result and the finite certificate needed by the next
lemma; the proof of totality uses the exact bound from the article.

\item \textbf{The Roth fragment.}
Formalize rational reduction, the rational-root divisibility theorem, the
finite no-root test, polynomial differences, and bisection.  Prove the finite
separation inequality directly at precision $k$.  The cutoff then follows
by exponent arithmetic.  The object statement is conditional only on the
finite input checks and a positive integer $a-(d-2)b$.  Do not declare a
general effective-Roth axiom to fill the missing higher-degree result.

\item \textbf{St\o rmer.}
Build the exact ordered quadratic-pair arithmetic.  Prove both pigeonhole
constructions for Pell, the quotient's integrality and positivity, and its
explicit bound.  Prove bounded generation by the fundamental solution.
Formalize the binomial identity and its three prime-index cases.  Finish
with the coefficient enumeration, smoothness filter, completeness, cardinality,
and height bounds.  No external finiteness theorem is required.

\item \textbf{The finite Davenport--Schmidt core.}
Start with a rational working point and a finite local-separation predicate
through height $2B^2$.  Formalize the pigeonhole estimate, exact rational
minimal records, integral-plane lemma, independence bound, determinant
identity, and finite growth contradiction.  Prove the existence of a triple
inside the explicit box.  The two-by-two minor coordinate reduction and
the six-term determinant expansion are finite polynomial calculations.

\item \textbf{Root certificates and the represented-real schema.}
For a linear or quadratic polynomial, define its finite root code, reduce a
rational root, and prove the height comparison and distance estimate.  Add
the single-sample precision lemma and the error transfer.  State this as a
schema for each fixed \texttt{PR} name and separation modulus, together with
their arithmetic correctness premises.  Add reflection and reciprocal
normalization.  The algebraic specialization needs the finite Sylvester
adjugate identity and its coefficient bound.

\item \textbf{Assemble and transport.}
Construct the final \texttt{Prf PRAAx} derivations.  Apply
\texttt{pra\_theorem\_iff} to transfer them to the concise presentation.
Check the resulting declarations and their dependencies in Lean.  A
metatheoretic theorem about \texttt{Rat}, \texttt{Int}, or \texttt{Real} is
only an intermediate validation milestone.
\end{enumerate}

The finiteness bounds in the positive proofs are intentionally generous.
Their role is to make every primitive recursion and every induction invariant
explicit.  Optimizing them is a separate project and is not needed for the
claimed proof-theoretic result.

\subsection{Exact computations supplied with the report}

The Python reference code uses only integers, fractions, and integer square
roots.  It has no third-party dependency.  The verification run performs the
following genuinely different checks.

First, the St\o rmer enumerator is compared with independent multiplicative
generation of smooth integers up to $100000$, for the empty prime set and
the sets $\{3,5\}$, $\{2\}$, $\{2,3\}$, and $\{2,3,5\}$.  The complete
enumeration theorem is proved above; the finite comparison checks that the
implementation produces the expected values.  Pell identities and the key
smooth-coordinate property are also checked for every nonsquare
$2\leq D\leq50$, at the first ten powers of the fundamental solution.

For $S=\{2,3,5\}$, the result is exactly
\[
 1,2,3,4,5,8,9,15,24,80
\]
as the smaller entries of the neighbor pairs.  The JSON evidence lists every
Pell coefficient and every retained or rejected candidate.  The implementation
does not allocate $B(D)$: it decides $y<\mathrm{base}^{\mathrm{exponent}}$
by a bounded number of divisions of $y$ by the base.  This is an exact
comparison, not a smaller heuristic search cutoff.

Second, the quadratic Roth example uses an independent exact decision
procedure in $\mathbb Q(\sqrt2)$.  For $\delta=1/2$, the proved cutoff is
$q<36$.  The independent check through $q=72$ finds the reduced exceptional
fractions
\[
 1,\quad 2,\quad \frac32,\quad\frac75.
\]
Each has a separate rational bisection certificate.  Because the theorem
excludes every $q\geq36$ and the numerator range is exhaustive, these are
all the exceptional reduced fractions for this example.  The larger
computational range is a regression check, not the reason completeness holds.

Third, bounded searches find the following real quadratic approximants to
$\sqrt[3]{2}$.  The displayed decimals are explanatory rounded values; the
JSON file stores the exact rational isolating intervals and exact rational
upper bounds.

\begin{center}
\begin{tabular}{@{}rrrr@{}}
\toprule
Minimal polynomial & Height & Error upper bound & $H^3$ times error\\
\midrule
$6X^2-2X-7$ & 7 & $3.480\cdot10^{-4}$ & $0.120$\\
$13X^2-14X-3$ & 14 & $1.430\cdot10^{-4}$ & $0.393$\\
$27X^2-11X-29$ & 29 & $1.222\cdot10^{-5}$ & $0.298$\\
$59X^2-41X-42$ & 59 & $9.379\cdot10^{-7}$ & $0.193$\\
\bottomrule
\end{tabular}
\end{center}

Each polynomial is primitive with positive nonsquare discriminant; its
reported height is therefore its actual algebraic height.  Opposite rational
endpoint signs certify its real root, and a separate cubic interval certifies
the target.  These examples illustrate the exponent-three phenomenon; they
do not execute the enormous general height bound.

For this target, the algebraic modulus construction is concrete.  With
$F(X)=X^3-2$, $d=3$, $A=2$, and $R=2$, it gives
\[
 \rho(H)=648\max(2,H)^4.
\]
The polynomial is irreducible by the elementary Eisenstein argument at $2$,
and bisection on $[1,2]$ supplies its real root name.  Thus the theorem applies
with explicit arithmetic data, despite the intentionally impractical general
bound.

\subsection{Reproduction and honest completion criteria}

Run \texttt{python3 code/verify.py} from the archive directory to regenerate
the exact computational evidence.  Run \texttt{./build.sh} to compile the
LaTeX article twice.  The build uses ordinary \texttt{pdflatex} and standard
LaTeX packages.  The archived PDF has been compiled and visually inspected.

The file \texttt{lean/FORMALIZATION\_STATUS.txt} lists the exact unexecuted
Lean build commands and the pending object-theory work.  No completion claim
for a machine-checked iPRA proof is made.  Likewise, no mathematical claim is
made that full effective Roth has been resolved.  A future completion of
that part would need either a specified weaker constructive conclusion with
its full arithmetic proof, or a new effective-height theorem.  Merely
rephrasing classical finiteness or importing it as a hypothesis would not
complete that task.

\clearpage
\begin{thebibliography}{99}

\bibitem{DavenportSchmidt1967}
H.~Davenport and W.~M.~Schmidt,
\emph{Approximation to real numbers by quadratic \mbox{irrationals}},
Acta Arithmetica \textbf{13} (1967), 169--176.
\href{https://doi.org/10.4064/aa-13-2-169-176}{doi:10.4064/aa-13-2-169-176}.
Primary article consulted, including its minimal-point and determinant argument.

\bibitem{Lehmer1964}
D.~H.~Lehmer, \emph{On a problem of St\o rmer},
Illinois Journal of Mathematics \textbf{8} (1964), 57--79.
\href{https://doi.org/10.1215/ijm/1256067456}{doi:10.1215/ijm/1256067456}.
The introduction describes St\o rmer's original coefficient method and
smooth-coordinate lemma.  The sharper Lehmer index bound is not used here.

\bibitem{Stormer1897}
C.~St\o rmer, \emph{Quelques th\'eor\`emes sur l'\'equation de Pell
$x^2-Dy^2=\pm1$ et leurs applications},
Videnskabsselskabets Skrifter, I. Matematisk-Naturvidenskabelig Klasse,
Christiania, no.~2 (1897), 48 pp.  Historical original.

\bibitem{Klazar2010}
M.~Klazar, \emph{St\o rmer's solution of the unit equation $x-y=1$},
author's exposition, 25 August 2010.
\url{https://kam.mff.cuni.cz/~klazar/stormer.pdf}.
Consulted as an exposition of the elementary prime-index proof.

\bibitem{Luckhardt1989}
H.~Luckhardt, \emph{Herbrand-Analysen zweier Beweise des Satzes von Roth:
Polynomiale Anzahlschranken},
Journal of Symbolic Logic \textbf{54} (1989), 234--263.
\href{https://doi.org/10.2307/2275028}{doi:10.2307/2275028}.

\bibitem{KohlenbachHerbrand}
U.~Kohlenbach, \emph{Herbrand's theorem and extractive proof theory},
author's manuscript, 1 September 2008, especially Theorem~1.9.
\url{https://www2.mathematik.tu-darmstadt.de/~kohlenbach/Gazette.pdf}.

\bibitem{AdiceamShirandami2026}
F.~Adiceam and V.~Shirandami,
\emph{Probabilistic effectivity in the Subspace Theorem},
Research in Number Theory \textbf{12}, article 8 (2026).
Published 13 January 2026.
\href{https://doi.org/10.1007/s40993-025-00692-0}{doi:10.1007/s40993-025-00692-0}.
The introduction describes the outstanding effective-height problem;
its probabilistic results are not a uniform cutoff for individual inputs.

\bibitem{ProveItPR}
V.~Reshetnikov and contributors, ProveIt,
\texttt{ConciseIPRA/PR.lean}, inspected at commit
\texttt{b33ee1acb02b7619dfef419e61c3e343735a7bc4},
15 November 2026.
\href{https://github.com/VladimirReshetnikov/ProveIt/blob/b33ee1acb02b7619dfef419e61c3e343735a7bc4/Logic/PrimitiveRecursiveArithmetic/ConciseIPRA/Lean/ConciseIPRA/PR.lean}{Pinned source: PR syntax, terms, axioms, and quantifier-free induction}.

\bibitem{ProveItDeduction}
V.~Reshetnikov and contributors, ProveIt,
\texttt{ConciseIPRA/Deduction.lean}, same commit.
\href{https://github.com/VladimirReshetnikov/ProveIt/blob/b33ee1acb02b7619dfef419e61c3e343735a7bc4/Logic/PrimitiveRecursiveArithmetic/ConciseIPRA/Lean/ConciseIPRA/Deduction.lean}{Pinned source: intuitionistic natural deduction and equality}.

\bibitem{ProveItBi}
V.~Reshetnikov and contributors, ProveIt,
\texttt{ConciseIPRA/BiInterpretation.lean}, same commit.
\href{https://github.com/VladimirReshetnikov/ProveIt/blob/b33ee1acb02b7619dfef419e61c3e343735a7bc4/Logic/PrimitiveRecursiveArithmetic/ConciseIPRA/Lean/ConciseIPRA/BiInterpretation.lean}{Pinned source: bi-interpretation and theorem transport}.

\end{thebibliography}

\end{document}
